\InputIfFileExists{../data/global.src}\relax\relax

\iffull
\definecolor{@joda}{RGB}{204, 186, 157}
\definecolor{@joda@}{RGB}{78, 68, 66}
\setbox\pinguA=\hbox{\tikz\pingu[body=pingu@green!52!pingu@black,body front=pingu@green!52!pingu@black!5!pingu@white,cloak=@joda,cloak cap=@joda,lightsaber right=pingu@green,lightsaber right length=1.4cm,shirt=@joda@,shirt no buttons,small,eyes wink,cane left=pingu@bronze!80!pingu@black,cane left length=11mm,cane left raise=1mm,right item angle=-30];}
\titlesuffix{\tikzpicture[overlay,remember picture]
    \node[above left,yshift=\btdmfootheight-1mm] at(current page.south east) {\copy\pinguA};
\endtikzpicture}
\def\packtitlenode#1#2{\node[font=\usebeamerfont{title},inner xsep=0pt,opacity=#1,outer ysep=2pt] (#2) {getauscht};}
\title[Erstes Tutorium -- Übungsblatt 2]{Diese Worte \texorpdfstring{\,\setbox0=\hbox{getau\kern5pt\relax scht}\phantom{\copy0}\llap{\tikz[baseline=(k.base)]{%
\packtitlenode0k
\def\kA{310}\def\kB{130}\def\kS{-.5\wd0-5pt}
\scope[xshift=-\kS]\packtitlenode0x
\clip (x.south west)--(x.\kA)--(x.\kB)-| cycle;
\packtitlenode1x
\endscope
\scope[xshift=\kS]
\packtitlenode0y
\clip (y.south east)--(y.\kA)--(y.\kB)-| cycle;
\packtitlenode1y
\endscope
% cutline
\fill[btdm@border@down,rounded corners=1.5pt] ([yshift=4pt,xshift=1.5pt]y.\kA)--([yshift=-4pt]y.\kB)--++(-3pt,0)--([xshift=-1.5pt,yshift=4pt]y.\kA) -- cycle;
\fill[btdm@border@down,rounded corners=1.5pt] ([yshift=4pt,xshift=-1.5pt]x.\kA)--([yshift=-4pt,xshift=-1.5pt]x.\kB)--++(3pt,0)--([xshift=1.5pt,yshift=4pt]x.\kA) -- cycle;
}\kern-.5\wd0}\;~}{getauscht} ich habe}
\subtitle{Tutorium Zwei}
\date{KW 19}
\addbibresource{references.bib}
\fi
\SetTutoriumNumber{2}

\iffull\begin{document}
\titleframe

\TopicOverview{2}
\fi

\iffull{\SummaryFrame
\def\sub#1#2{\node[font=\tiny\sffamily,align=center,gray,below=-3mm] at(#1.south) {\strut#2\strut};}
\begin{frame}[fragile,c]{Kurzwiederholung}
    \begin{itemize}[<+(1)->]
        \itemsep14pt
        \item Algorithmenkonstruktion\smallskip\par \centerline{%
            \begin{tikzpicture}
                \onslide<+(1)->{\node (0) at(0,0) {\strut Spezifikation};
                \sub0{Begriffe mit\\Problemrelevanz}}
                \foreach[count=\i,remember =\i as \li (initially 0)] \a/\t in {Abstraktion/{Gegeben \& Gesucht},{\makebox[14mm]{\only<8->{\sbseries}Entwurf}}/{Algorithmus},Verifikation/{Termination \&\\partielle Korrektheit},Aufwandsanalyse/{Laufzeitverhalten}} {
                    \onslide<+(1)->{\node[right=.5mm] (k\i) at(\li.east) {\strut\faAngleRight};
                    \node[right=.5mm] (\i) at (k\i.east) {\strut\a};
                    \sub\i{\t}}
                }
            \end{tikzpicture}}\vspace*{-\medskipamount}
        \item<9-> Kerneigenschaften eines Algorithmus: \begin{itemize}
            \item<10-> Von einem Prozessor \info{Mensch,~\ldots} in endlich vielen Schritten Ausführ- und Reproduzierbar
            \item<11-> Eine endliche Beschreibung in Elementaroperationen.
        \end{itemize}
        \item<12-> Abstrakte Variablenoperationen:
\begin{plainjava}
!*\onslide<13->*!long x; !*\tikzmarknode{@1}{\phantom{Ig}}*!
!*\onslide<13->*!x = 21; !*\tikzmarknode{@2}{\phantom{Ig}}*!
!*\onslide<13->*!x = 42; !*\tikzmarknode{@3}{\phantom{Ig}}*!
\end{plainjava}
    \end{itemize}
\begin{tikzpicture}[overlay,remember picture,T/.style={sol@colors@lst@comments,font=\itshape\sffamily\footnotesize}]
\onslide<14->{\node[right,T] at(@1.east) {\textbf{Deklaration} --- Reservieren von Namen mit Typ long};}
\onslide<15->{\node[right,T] at(@2.east) {\textbf{Initialisierung} --- Erste Wertzuweisung, macht die Variable verwendbar};}
\onslide<16->{\node[right,T] at(@3.east) {\textbf{Wertzuweisung} --- Überschreiben/Ändern des vorherigen Wertes};}
\end{tikzpicture}
\onslide<1->
\end{frame}
}\fi

\SetNextSectionText{\say{Machs farbig}\\--- Isabell Bannweg}
\section{Exkurs: Kommandozeilenparameter}
\begin{frame}[fragile]{M-Array me und andere Zukunftsträume}
    \begin{itemize}[<+(1)->]
        \itemsep9pt
        \item Arrays werden später genauer behandelt.
        \item Arrays sind Listen fester Größe die Elemente des gleichen Datentyps beinhalten.
        \item Wir notieren sie mit \T{[]}:
\begin{plainjava}[lineskip=2pt]
!*\onslide<5->*!int!*\tikzmarknode{@}{\sbasic{\HStrut[]}}*! arr = {!*\tikzmarknode{@0}{\snum{\HStrut42}}*!, !*\tikzmarknode{@2}{\snum{\HStrut-7}}*!,  !*\tikzmarknode{@4}{\snum{\HStrut13}}*!};

!*\onslide<7->*!System.out.println(!*\tikzmarknode{@3}{\sbasic{arr[\snum{\HStrut1}]}}*!);!*\onslide<9->*! // :yields: -7

!*\onslide<10->*!System.out.println(!*\tikzmarknode{@5}{\sbasic{\HStrut arr.length}}*!);!*\onslide<12->*! // :yields: 3

!*\onslide<13->*!System.out.println(!*\tikzmarknode{@6}{\sbasic{arr[arr.length - \snum{\HStrut1}]}}*!); !*\onslide<15->*!// :yields: 13
!*\onslide<1->*!
\end{plainjava}
    \end{itemize}
\begin{tikzpicture}[overlay,remember picture,codeouthl]
    \def\hlhs{0mm}
    \onslide<6->{
        \hlcode[-.35pt]{@}{array}
        \draw[Kite-] ([yshift=-.2mm]@.south) to[out=290,in=180] ++(1,-.35) node[right] {\small\itshape Ein Array an Integern};
    }
    \onslide<8->{
        \def\hlopa{.6}
        \def\hlcolor{pingu@yellow}\hlcode[-.35pt]{@2}{-7}
        \hlcode[-.35pt]{@3}{1}
        \node at(@2) {\snum{\HStrut-7}};
        \node at(@3) {\sbasic{\HStrut arr[\snum{1}]}};
    }
    \onslide<9->{\draw[Kite-] ([yshift=-.2mm]@3.330) to[out=240,in=0] ++(-1,-.25) node[left] {Beginnt bei \(0\)};}
    \onslide<14->{
        \def\hlopa{.6}
        \def\hlcolor{pingu@green}\hlcode[-.35pt]{@4}{13}
        \hlcode[-.35pt]{@6}{1}
        \node at(@4) {\snum{\HStrut13}};
        \node at(@6) {\sbasic{\HStrut arr[arr.length - \snum{1}]}};
    }
    \onslide<11->{
        \def\hlopa{.6}\def\hlcolor{pingu@red}
        \hlcode[-.35pt]{@5}{1}
        \node at(@5) {\sbasic{\HStrut arr.length}};
        \node[above=-1.25mm,opacity=.6,pingu@red] at(@0.north) {\textbullet};
        \node[above=-1.25mm,opacity=.6,pingu@red] at(@2.north) {\textbullet};
        \node[above=-1.25mm,opacity=.6,pingu@red] at(@4.north) {\textbullet};
    }
    \onslide<12->{\draw[Kite-] ([yshift=-.2mm]@5.335) to[out=240,in=0] ++(-1,-.3) node[left] {Liefert die Array-Größe};}
\end{tikzpicture}
\end{frame}

\begin{frame}[fragile]{Die Sache mit dem \T{\bfseries String[] args}}
% I am not using \SetupLstHl because i want the fade out... of course i do... :c
\begin{uncoverenv}<2->
\begin{plainjava}[lineskip=4pt]
!*\only<2|handout:0>{\sbasic{\skA{public class} Example \{}}\only<3->{\sbasic{\color{codeouthl}public class Example \{}}*!
    !*\only<2|handout:0>{\sbasic{\skA{public} \skB{static void} main(}}\only<3->{\sbasic{\color{codeouthl}public static void main(}}*!String[] args!*\only<2|handout:0>{\sbasic{) \{}}\only<3->{\sbasic{\color{codeouthl}) \{}}*!
        !*\onslide<6->*!System.out.println(args.length);!*\onslide<1->*!
    !*\only<2|handout:0>{\sbasic{\}}}\only<3->{\sbasic{\color{codeouthl}\}}}*!
!*\only<2|handout:0>{\sbasic{\}}}\only<3->{\sbasic{\color{codeouthl}\}}}*!
\end{plainjava}
\end{uncoverenv}
\begin{itemize}[<+(1)->]
    \itemsep6pt
    \item<4-> Hier haben wir ein Array von Strings.
    \item<5-> Dieses wird von der Java Virtual Machine gefüllt.
    \item<7-> Beim Programmstart mit \cbash{java Example}, können wir Argumente übergeben.
\end{itemize}
\end{frame}

\iffull
\MakeThePinguExplainIt[text width=3.25cm]{cap=!hide,eye patch right,eye patch left,headband=paletteA,bill=flat,lollipop left,right item angle=-30,cup=!hide,glasses=!hide,body type=legacy}{Fürs erste werden wir es dabei belassen.}
\begin{frame}[fragile,c]{Ein wunderbares Beispiel}
\SetupLstHl
\begin{plainjava}[lineskip=4.5pt]
!*\onslide<2->*!|ihl|public class Example {|ihl|
!*\onslide<2->*!    |ihl|public static void main(|ihl|String[] !*\tikzmarknode{args}{\sbasic{\HStrut args}}*!|ihl|) {|ihl|

!*\onslide<3->*!        for(int !*\tikzmarknode{i}{\sbasic{\HStrut i}}*! = 0; !*\tikzmarknode{i4}{\sbasic{\HStrut i}}*! < !*\tikzmarknode{args-length}{\sbasic{\HStrut args.length}}*!; !*\tikzmarknode{i2}{\sbasic{\HStrut i}}*!++) {
!*\onslide<4->*!            System.out.println(!*\tikzmarknode{args-a}{\sbasic{\HStrut args[}\kern-1pt}\kern1pt\tikzmarknode{i3}{\sbasic{\HStrut i}}\kern1pt\tikzmarknode{args-b}{\kern-1pt\sbasic{\HStrut ]}}*!);
!*\onslide<3->*!        }

!*\onslide<2->*!    |ihl|}|ihl|
!*\onslide<2->*!|ihl|}|ihl|
\end{plainjava}
\begin{tikzpicture}[overlay,remember picture,T/.style={sol@colors@lst@comments,font=\itshape\sffamily\footnotesize}]
    \onslide<5->{\node[below left,xshift=-5mm,yshift=-1.65cm,outer sep=0pt, inner sep=0pt] (@) at(current page.north east) {\bbash{Hermeline mögen Mäuse}};
    \node[left=0pt,outer sep=0pt, inner sep=0pt] (@2) at(@.west) {\bbash{java Example}\T{~}\null};}
    \def\hlhs{.1mm}
    \onslide<6->{
        \draw[decoration={brace,aspect=.85},decorate] (@.south east) to coordinate[pos=.85] (@3) (@.south west);
        \draw[-Kite] ([yshift=-4pt]@3.south) to[out=220,in=40] ([yshift=3pt]args.north);
    }
    \onslide<7->{
        \def\hlopa{.5}
        \def\hlcolor{pingu@blue}\hlcode[-.35pt]{args}{@args}
        \node at(args) {\sbasic{\HStrut args}};
        \node[T,below right] at(args.south) {args = \{\T{\dq}Hermeline\T{\dq}, \T{\dq}mögen\T{\dq}, \T{\dq}Mäuse\T{\dq}\}};
    }
    \onslide<8->{
        \def\hlopa{.5}
        \def\hlcolor{pingu@blue}\hlcode[-.35pt]{args-length}{@args}
        \node at(args-length) {\sbasic{\HStrut args.length}};
        \node[above,pingu@blue] at (args-length.north) {\sbseries\itshape 3};
    }
    \onslide<10->{{
        \def\hlopa{.5}
        \def\hlcolor{pingu@blue}
        \hlcode[-.35pt]{args-a}{@args}\node at(args-a) {\sbasic{\HStrut args[\kern-1pt}};
        \hlcode[-.35pt]{args-b}{@args}\node at(args-b) {\sbasic{\kern-1pt\HStrut ]}};
    }}
    \onslide<9->{
        \def\hlopa{.275}
        \def\hlcolor{pingu@yellow}
        \hlcode[-.35pt]{i}{@i}\node at(i) {\sbasic{\HStrut i}};
        \hlcode[-.35pt]{i2}{@i}\node at(i2) {\sbasic{\HStrut i}};
        \hlcode[-.35pt]{i3}{@i}\node at(i3) {\sbasic{\HStrut i}};
        \hlcode[-.35pt]{i4}{@i}\node at(i4) {\sbasic{\HStrut i}};
    }
    \onslide<11->{ \node[below right,yshift=-1mm] (@a) at (i3.south) {\T{Hermeline}}; \draw[very thick,pingu@blue!50!btdm@background,-Kite,rounded corners=2pt] (args-a.south) |- (@a); \node[left=1cm,pingu@yellow] at (@a.west) {\T{\sbseries i\,=\,0}};}
    \onslide<12->{ \node[below right,yshift=-.635cm] (@b) at (i3.south) {\T{mögen}}; \draw[very thick,pingu@blue!50!btdm@background,-Kite,rounded corners=2pt] (args-a.south) |- (@b); \node[left=1cm,pingu@yellow] at (@b.west) {\T{\sbseries i\,=\,1}}; }
    \onslide<13->{ \node[below right,yshift=-1.2cm] (@c) at (i3.south) {\T{Mäuse}}; \draw[very thick,pingu@blue!50!btdm@background,-Kite,rounded corners=2pt] (args-a.south) |- (@c); \node[left=1cm,pingu@yellow] at (@c.west) {\T{\sbseries i\,=\,2}};}
    \onslide<0|handout:1>{\node[left=-2mm,xshift=5mm,scale=.8,yshift=\btdmfootheight] at(current page.-16) {\copy\pinguexplainbox};}% copy for animations
\end{tikzpicture}
\end{frame}
\fi

%     \item<8-> Mit \bbash{java Example Hermeline mögen Mäuse} ist \bjava{args = \{"Hermeline", "mögen", "Mäuse"\}}.

\begin{frame}
    TODO: Wrapper-klassen für Integer-parseInt und so
\end{frame}

\SetNextSectionText{Es gibt keine if-Schleifen, sondern nur if-Abfragen!\\--- \href{http://if-schleife.de/}{if-schleife.de}}
\section{Präsenzaufgabe}
\begin{frame}[fragile,c]{Präsenzaufgabe}
\begin{aufgabe}{Wenn\ldots\ Ja wenn nur die Schleife nicht wär\ldots}
    \onslide<2->{Legen Sie eine Java Datei namens \T{ErsteSchleife.java} an \info{oder bearbeiten Sie die Aufgabe auf einem Blatt Papier}.} \onslide<3->{Lesen Sie in der main Methode eine \bjava{int} Variable namens \(n\) über die Kommandozeilenparameter ein und implementieren Sie eine for Schleife, die jede dritte Zahl von \(1\) bis einschließlich \(n\) ausgibt.}
    \onslide<4->Beispiel:
\begin{plainjava}
n = 13 !*\onslide<+(1)->*!// :yields: 1 4 7 10 13
\end{plainjava}
\onslide<1->
\end{aufgabe}
\end{frame}

\begin{frame}{}

\end{frame}

\section{Abschließendes}
{\SummaryFrame
\begin{frame}[t]{Zusammenfassend}
\pause \printBibCommand
\vfill\vfill % double fill for more fraction
\begin{itemize}[<+(1)->]
    \itemsep14pt
    \item Wir haben die Geburt von Variablen sowie ihre Wachstumsphase kennengelernt: \begin{itemize}
        \item Die \textit{Deklaration} (\bjava{int x}) reserviert einen Namen samt Charakteristika (Typ,~\ldots)
        \item Die \textit{Initialisierung} (\bjava{int x; x = 5}) beschreibt die erste Wertzuweisung
        \item Eine \textit{(Wert-)Zuweisung} (\bjava{x = 3}) ist jede weitere Änderung des gespeicherten Wertes
    \end{itemize}
    \item In die Wahl des \say{richtigen\texttrademark} Datentyps fließen viele Informationen mit ein.
\end{itemize}
\end{frame}
}

\outro{\vskip6mm\centering\begin{tikzpicture}[scale=2.5]
    % \only<2->{\pingu[monocle left,right eye wink,left eye vertical,laptop right,cane left,tie,body type=legacy,headphones=pingu@green!80!pingu@black]}
    % TODO
\end{tikzpicture}}


\iffull\end{document}\fi
